Rhythm games have developed distinctive techniques and cultures over their long history. Countless players have accumulated knowledge and skills up to the present day. These activities continue to deepen and expand, and demand for scholarly and cultural inquiry into rhythm games continues to grow.

Rhythm Games Academy was established as an academic fan community that brings together knowledge, techniques, and cultures relating to rhythm games, conscious of its mission to carry forward and develop this rich body of work. The Academy pledges to support its members' independent activities, including research, creation, courses, and reflection, and to contribute to the development of rhythm game culture as a whole through free and open exchange and inquiry.

To fulfill these missions together with its members, the Academy hereby adopts this Charter.

\begin{articles}
  \item Rhythm Games Academy is an academic fan community whose permanent mission is to gather, pass on, and develop all intellectual and cultural heritage relating to rhythm games. It respects each member's inquiry and creativity, and regards sharing the results with rhythm game culture as a whole as its reason for existing.
  \item On the basis of academic freedom, Rhythm Games Academy respects members' free and independent activities as its guiding principle. Regardless of the form or subject of activities, including research, creation, courses, and reflection, it supports equally the inquiry of everyone who wishes to engage independently with rhythm games. It believes that the exchange of diverse perspectives and interests produces richer knowledge.
  \item Rhythm Games Academy must respect all members as individuals and treat them equally. Skill level, years of experience, and achievements do not establish a personal hierarchy among members. Members engage with one another as equals from their respective positions and respect one another's knowledge and activities.
  \begin{paragraphs}
    \item Members have a duty to respect one another's freedom and maintain healthy order within the community.
  \end{paragraphs}
  \item Rhythm Games Academy opens its doors to everyone who wishes to engage independently with rhythm games and can comply with the Academy's rules. Participation in activities within the Academy is a matter of each member's free choice, and all forms of participation are equally welcome.
  \item Rhythm Games Academy provides places and opportunities for members to pursue inquiry on a continuing basis. It maintains an environment for presenting, discussing, and preserving their work, and supports this foundation regardless of the field of inquiry or form of activity. This environment is a shared resource of all members, and the Academy must continually endeavor to improve it.
  \item Rhythm Games Academy regards it as an important responsibility to preserve the results of members' activities on a continuing basis, with due regard for members' rights, and to pass them on to future generations as part of rhythm game culture's heritage. Through members' courses and research activities, it passes the knowledge of earlier generations to the next and maintains continuity in scholarly activity within the Academy.
  \item Rhythm Games Academy seeks to sustain its activities over the long term, beyond any temporary burst of enthusiasm. It maintains accumulated knowledge and organizational autonomy so that its mission can be carried forward despite changes in membership and the times. It aims to become a lasting cultural community within the rhythm game community.
  \item The Management of Rhythm Games Academy consists of the Director and other officeholders specified in the Academy's rules. Management bears responsibility for administering the Academy and manages its affairs in accordance with the principles of this Charter. Those who constitute Management are themselves members and are bound by this Charter. The details of how Management is selected and of its duties are provided in separately established rules.
  \item Any Academy Regulations or acts of Management that conflict with this Charter have no effect to the extent of that conflict, whether in whole or in part.
  \begin{paragraphs}
    \item All members have the right to raise an objection to a violation of this Charter by Management.
  \end{paragraphs}
  \item An amendment to this Charter may be proposed by Management or by a joint petition signed by at least one fifth of all members. An amendment requires the support of at least two thirds of the members who participate in the vote and of at least one fifth of all members. Management must confirm that these requirements have been met before putting the amendment into effect. Members must be informed of amendments, and efforts must be made to reflect a broad range of views.
\end{articles}

\begin{samepage}
\bigskip
\noindent Supplementary provision\\
This Charter takes effect on 1 August 2026.

\bigskip
\noindent History\\
Adopted on 22 June 2026.
\end{samepage}
